% R13 THEORY PREWORK -- NEW FILE ONLY; not included by main.tex or supplement.tex.
% Base: branch scgi-ceiling-diagnostic-r1, HEAD bcbb31254c343eb04e8cb96d7c3157bfc9818b36.
% Purpose: paste-ready material for the gauge, direct Poisson lower bound,
% oracle Fisher reference, and permutation-scope repairs.

\paragraph{Proposition (exact collision and quotient tangent).}
Let
\[
 \mu(\ell,T)=D_{e^\ell}MT,
 \qquad \ell\in S\le\mathbb R^N,
 \qquad T\in\mathbb R^K,
\]
where $S$ is a known linear subspace, $\mathbf1\in S$, and suppose
$c=MT$ has no zero entries.  For
any $u\in S$, the pair $(\ell+u,T_u)$ produces the same noiseless
record as $(\ell,T)$ if and only if
\begin{equation}
 MT_u=D_{e^{-u}}MT. \tag{G.1-R13}
\end{equation}
Consequently, in the square case $N=K$, if
$M\in\mathbb R^{K\times K}$ is invertible, every
$u\in S$ generates the exact collision
\begin{equation}
 \ell_u=\ell+u,\qquad
 T_u=M^{-1}D_{e^{-u}}MT. \tag{G.2-R13}
\end{equation}
If $M\in\mathbb R^{N\times K}$ is tall with full column rank, such a
collision exists if and only if
\begin{equation}
 D_{e^{-u}}MT\in\operatorname{range}(M), \tag{G.3-R13}
\end{equation}
equivalently $(I-P_M)D_{e^{-u}}MT=0$.  The global-scale gauge is the
special case $u=t\mathbf1$, for which $T_u=e^{-t}T$.  Under the
nonvanishing-carrier assumption, $T_u$ is proportional to $T$ only
when $u$ is constant.

The differential of the mean map in a direction
$(u,h)\in S\times\mathbb R^K$ is
\[
 D\mu_{(\ell,T)}[u,h]
 =D_{e^\ell}\{D_{MT}u+Mh\}.
\]
Its null space is
\[
 \mathcal N_{\ell,T}
 =\{(u,h):D_{MT}u+Mh=0\},
\]
and always contains the gauge tangent
$\operatorname{span}\{(\mathbf1,-T)\}$.  The model is locally
differentially identifiable modulo scale if and only if
\begin{equation}
 \mathcal N_{\ell,T}
 =\operatorname{span}\{(\mathbf1,-T)\}. \tag{G.4-R13}
\end{equation}

\paragraph{Proof.}
Equality of records is
$D_{e^{\ell+u}}MT_u=D_{e^\ell}MT$; multiplication by
$D_{e^{-\ell-u}}$ gives (G.1-R13).  Square inversion gives
(G.2-R13), while tall solvability is exactly (G.3-R13).  If
$T_u=\gamma T$, then $D_{e^{-u}}c=\gamma c$; because every $c_n\ne0$,
$e^{-u_n}=\gamma$ for every $n$, so $u$ is constant.  Differentiation
at zero gives the tangent formula and (G.4-R13).
\hfill$\blacksquare$

\paragraph{Quotient Fisher consequence.}
Let $J$ be the Jacobian above.  In a regular Poisson model with
strictly positive mean,
\[
 I=J^\top D_\mu^{-1}J,
\]
and in a Gaussian model with known positive-definite covariance
$\Sigma$,
\[
 I=J^\top\Sigma^{-1}J.
\]
The weighting matrix is positive definite, hence
\[
 \ker I=\ker J=\mathcal N_{\ell,T}.
\]
The Fisher metric is positive definite on the quotient tangent space
precisely when (G.4-R13) holds.  An exact collision curve implies a
Fisher-null direction.  The converse yields only differential
nonidentifiability unless an appropriate constant-rank/fibre condition
or an explicit construction shows that the null direction integrates
to an exact collision curve.

The collision statement is on the unconstrained object domain
$\mathbb R^K$.  For a nonnegative object class it gives a local
nonidentifiability statement only at strictly positive interior objects
and for sufficiently small perturbations; it does not automatically
preserve boundary objects with zero pixels.

\paragraph{Theorem (oracle Poisson pointwise lower bound).}
Let $0<\beta\le1$ and, conditionally on known positive carriers
$q_1,\ldots,q_N$, observe independent counts
\[
 C_i\sim\operatorname{Pois}(q_i e^{\ell_i}),
 \qquad i=1,\ldots,N.
\]
Let $P_N=I-N^{-1}\mathbf1\mathbf1^\top$, and estimate $P_N\ell$.
Fix an interior index $n_0$, a baseline $\ell^{(0)}$ lying strictly
inside both the amplitude and H\"older constraints (for example,
$\|\ell^{(0)}\|_\infty\le a_0/2$ and seminorm at most $L_a/2$), and a
window $A$ of order $W$
containing $n_0$.  Choose a contrast
$\psi^{(W)}$ satisfying
\[
 \operatorname{supp}\psi^{(W)}\subseteq A,\quad
 \mathbf1^\top\psi^{(W)}=0,quad
 \|\psi^{(W)}\|_\infty\le1,quad
 \psi^{(W)}_{n_0}=1,
\]
\[
 \|\psi^{(W)}\|_2^2\ge cW,\qquad
 |\psi^{(W)}_i-\psi^{(W)}_j|
 \le C W^{-\beta}|i-j|^\beta.
\]
Put $\lambda_i^{(0)}=q_i e^{\ell_i^{(0)}}$ and
\[
 I_A(\psi)=\sum_{i\in A}\lambda_i^{(0)}\psi_i^2
 \le W\bar\lambda_A,
 \qquad
 \bar\lambda_A={1\over W}\sum_{i\in A}\lambda_i^{(0)}.
\]
For
\[
 a\le c_0\min\{a_0,L_aW^\beta,I_A(\psi)^{-1/2}\},
\]
both $\ell^{(0)}$ and
$\ell^{(1)}=\ell^{(0)}+a\psi^{(W)}$ belong to a rescaled
$H_\beta(L_a)$ ball, and
\begin{equation}
 \inf_{\widehat\ell}\max_{r\in\{0,1\}}
 \mathbb E_r\!\left[
  \{(P_N\widehat\ell)_{n_0}-(P_N\ell^{(r)})_{n_0}\}^2
 \right]
 \ge ca^2. \tag{L.1-R13}
\end{equation}
In particular,
\begin{equation}
 R_{n_0}^{\mathrm{quot}}
 \ge c\min\left\{a_0^2,L_a^2W^{2\beta},
                   {1\over W\bar\lambda_A}\right\}. \tag{L.2-R13}
\end{equation}

\paragraph{Proof.}
For Poisson means
$\lambda_i^{(1)}=\lambda_i^{(0)}e^{a\psi_i}$,
\[
 \operatorname{KL}(P_0,P_1)
 =\sum_i\lambda_i^{(0)}
   \{e^{a\psi_i}-1-a\psi_i\}
 \le {e^{a_0}\over2}a^2I_A(\psi).
\]
The chosen amplitude keeps the divergence below a numerical constant.
Because $\mathbf1^\top\psi=0$, quotient projection does not reduce the
separation at $n_0$.  Pinsker followed by Le Cam's two-point lemma gives
(L.1-R13), and $I_A(\psi)\le W\bar\lambda_A$ gives (L.2-R13).
\hfill$\blacksquare$

When local photon rates have common order $\bar\lambda$, optimizing
$W$ yields
\[
 R_{n_0}^{\mathrm{quot}}
 \gtrsim L_a^{2/(2\beta+1)}
          \bar\lambda^{-2\beta/(2\beta+1)}.
\]
In normalized time, $L_a=LN^{-\beta}$, this is
\[
 R_{n_0}^{\mathrm{quot}}
 \gtrsim L^{2/(2\beta+1)}
 (N\bar\lambda)^{-2\beta/(2\beta+1)}.
\]
These optimized displays assume $1\lesssim W_*\lesssim N$; otherwise
the lower bound remains in the one-frame or finite-record saturation
regime already visible in (L.2-R13).

\paragraph{Theorem (integrated Poisson--Assouad lower bound).}
Assume the baseline lies strictly inside the amplitude and H\"older
ball, for example with seminorm at most $L_a/2$.  Partition an interior
fraction of $\{1,\ldots,N\}$ into $m\asymp N/W$ cells of length $CW$.
Inside each cell reserve a guard band of order $W$ and place a translated,
rescaled copy $\psi_j$ of one fixed compactly supported, mean-zero
$H_\beta$ template.  Require that $\psi_j$ vanish on the cell boundary,
that the active supports be separated by at least $cW$, and that
$\|\psi_j\|_2^2\ge cW$.  Define
\[
 \ell^\omega=\ell^{(0)}
 +a\sum_{j=1}^m\omega_j\psi_j,
 \qquad \omega\in\{0,1\}^m,
\]
where
\[
 \bar\lambda_*=\max_j |A_j|^{-1}
  \sum_{i\in A_j}\lambda_i^{(0)},
\]
and
\[
 a\le c_0\min\{a_0,L_aW^\beta,(W\bar\lambda_*)^{-1/2}\}.
\]
Then
\begin{equation}
 \inf_{\widehat\ell}\sup_\omega
 {1\over N}\mathbb E_\omega
 \|P_N(\widehat\ell-\ell^\omega)\|_2^2
 \ge ca^2, \tag{A.1-R13}
\end{equation}
and hence
\begin{equation}
 R_{\mathrm{int}}^{\mathrm{quot}}
 \ge c\min\left\{a_0^2,L_a^2W^{2\beta},
                   {1\over W\bar\lambda_*}\right\}. \tag{A.2-R13}
\end{equation}

\paragraph{Proof.}
Construct the fixed template from two smooth bumps with opposite signs.
Its discrete sum is zero, its squared norm after rescaling is order $W$,
and its H\"older seminorm is order $W^{-\beta}$.  Within one active
support the scaled-template bound applies.  Across distinct supports,
the guard-band separation is at least $cW$, while the perturbation
amplitude is at most $a\le cL_aW^\beta$; pairs involving a guard band
use the template's vanishing boundary.  Hence the H\"older constant of
the complete sum is bounded independently of $m$, and the interior
baseline assumption keeps every alternative in the parameter ball.
Adjacent hypercube
vertices differ on one block, and the preceding Poisson calculation
gives
\[
 \operatorname{KL}(P_\omega,P_{\omega^{(j)}})
 \le Ca^2W\bar\lambda_*\le\kappa<1.
\]
All pairwise differences have zero mean, so quotient projection loses
no distance:
\[
 \|P_N(\ell^\omega-\ell^{\omega'})\|_2^2
 \ge ca^2W d_H(\omega,\omega').
\]
Assouad's lemma then gives
\[
 {1\over N}\inf_{\widehat\ell}\sup_\omega
 \mathbb E_\omega\|P_N(\widehat\ell-\ell^\omega)\|_2^2
 \ge c{mW\over N}a^2\ge c'a^2.
\]
\hfill$\blacksquare$

If $q_i=\Lambda_0B_i$ is random with a parameter-independent law and
$\mathbb EB_i<\infty$, choose every alternative and its amplitude
before observing $B$, using the expected block information
\[
 \overline{\mathcal I}_A
 =\mathbb E_B\sum_{i\in A}\lambda_i^{(0)}\psi_i^2.
\]
The same calculation then applies to the carrier-observed oracle
experiment $(C,B)$ because its joint KL is
\[
 \operatorname{KL}(P_0^{C,B},P_1^{C,B})
 =\mathbb E_B\operatorname{KL}
   (P_0^{C\mid B},P_1^{C\mid B}).
\]
Marginalization gives
\[
 \operatorname{KL}(P_0^C,P_1^C)
 \le\operatorname{KL}(P_0^{C,B},P_1^{C,B}),
\]
so the oracle lower bound transfers to the less-informative count-only
experiment.  This transfers hardness; it does not construct a blind
estimator.  For dark count $d>0$, replace total count information by
\[
 \overline{\mathcal I}_A(\psi)
 =\mathbb E_B\sum_{i\in A}{s_i^2\over s_i+d}\psi_i^2,
 \qquad s_i=\Lambda_0e^{\ell_i}b_i.
\]
This is the local quadratic-KL information in the contrast direction;
the bounded-perturbation remainder must be controlled uniformly.

\paragraph{Proposition (oracle local-shift information).}
For a fixed window $A$, suppose independently
\[
 C_k\sim\operatorname{Pois}\{\lambda_k(t)\},
 \qquad \lambda_k(t)=s_ke^t+d_k,
\]
where $s_k>0$ and $d_k\ge0$ are known.  The score and Fisher information
for the one-dimensional common shift $t$ are
\[
 S_t(C)=\sum_{k\in A}(C_k-\lambda_k(t))
        {s_ke^t\over\lambda_k(t)},
\]
\begin{equation}
 I_A(t)=\sum_{k\in A}{s_k^2e^{2t}\over s_ke^t+d_k}. \tag{F.1-R13}
\end{equation}
When $d_k=0$, $I_A(t)=\sum_{k\in A}\lambda_k(t)$.  Therefore
$1/I_A(t)$ is the scalar CR reference for a regular unbiased estimator
of a common log-intensity shift with carriers, offsets, and all other
path coordinates held fixed.  It is not the nuisance-path quotient
CRB.

For bias $b(t)=\mathbb E_t\widehat t-t$, the biased CR inequality is
\[
 \operatorname{MSE}_t(\widehat t)
 \ge b(t)^2+{\{1+b'(t)\}^2\over I_A(t)}.
\]
Thus MSE below $1/I_A$ is compatible with shrinkage bias, while MSE
above $1/I_A$ does not establish unbiasedness.

For a path model $\ell=U\theta$ at $d=0$, where the columns of $U$
represent estimable contrast coordinates separated from the constant
direction, and with
$D_\lambda=\operatorname{diag}(\lambda_1,\ldots,\lambda_N)$,
\[
 I_\theta=U^\top D_\lambda U.
\]
If a common nuisance scale $m$ enters as $\ell+m\mathbf1$, eliminating
$m$ gives
\[
 I_{\theta\mid m}=U^\top D_\lambda U
 -U^\top D_\lambda\mathbf1
  (\mathbf1^\top D_\lambda\mathbf1)^{-1}
  \mathbf1^\top D_\lambda U.
\]
This matrix, or its pseudoinverse for an estimable contrast, is the
quotient Fisher object in this known-carrier parametric path model.  It
does not include unknown-object or unknown-carrier nuisance parameters;
a contrast outside the range of the information matrix has no finite CR
bound.  Overlapping windows also create
cross-window covariance, so an average of $1/I_n$ is not an exact
whole-record centered-loss CRB.

\paragraph{Proposition (two permutation semantics).}
Let $z=(z_1,\ldots,z_K)$ be fixed.  If $\Pi$ is uniform on all
permutations and $Z_n=z_{\Pi(n)}$, then $(Z_1,\ldots,Z_K)$ is
exchangeable over the randomization draw.  Writing
$\sigma_z^2=K^{-1}\sum_j(z_j-\bar z)^2$,
\[
 \operatorname{Var}(Z_i)=\sigma_z^2,
 \qquad
 \operatorname{Cov}(Z_i,Z_j)=-{\sigma_z^2\over K-1}\quad(i\ne j),
\]
and for a fixed set of $W$ time positions,
\[
 \operatorname{Var}\left({1\over W}\sum_{i\in A}Z_i\right)
 ={\sigma_z^2\over W}{K-W\over K-1}.
\]

Now let $G=\mathbb F_2^m$, $K=2^m$, let $\chi_g$ be Walsh characters,
and let
\[
 Z_g(P)=\sum_{j\in G}\chi_g(j)T_{Pj},
\]
where one uniform pixel permutation $P$ is shared by all rows.  For
every nonzero $g$, the one-dimensional distribution of $Z_g(P)$ is
the same, with
\[
 \mathbb E_PZ_g=0,
 \qquad
 \operatorname{Var}_PZ_g={KS_2-S_1^2\over K-1}.
\]
However, the vector $\{Z_g(P):g\ne0\}$ is not exchangeable in general.
Let $\bar T=K^{-1}\sum_jT_j$ and
$S_{3,c}=\sum_j(T_j-\bar T)^3$.  For distinct nonzero rows $g,h,r$,
\begin{equation}
 \mathbb E_P[Z_gZ_hZ_r]=
 \begin{cases}
 \displaystyle {K^2\over(K-1)(K-2)}S_{3,c},&g+h+r=0,\\[5pt]
 0,&g,h,r\text{ linearly independent}.
 \end{cases} \tag{P.1-R13}
\end{equation}
For $K\ge8$ and a generic object with $S_{3,c}\ne0$, relational and
independent triples therefore have different joint laws, disproving
full row exchangeability under a common pixel permutation.

\paragraph{Proof of (P.1-R13).}
Replace $T_j$ by $x_j=T_j-\bar T$ without changing non-DC Walsh
coefficients.  For the without-replacement sample $X_i=x_{P(i)}$,
\[
 \mathbb EX_i^3={S_{3,c}\over K},\qquad
 \mathbb EX_i^2X_j=-{S_{3,c}\over K(K-1)}\quad(i\ne j),
\]
\[
 \mathbb EX_iX_jX_k={2S_{3,c}\over K(K-1)(K-2)}
 \quad(i,j,k\text{ distinct}).
\]
If $g+h+r=0$, character orthogonality gives sign sums $K$, $-K$ for
each exactly-two-equal pairing, and $2K$ for the all-distinct class.
Substitution gives the first line of (P.1-R13).  If $g+h+r\ne0$, all
corresponding character sums vanish.
\hfill$\blacksquare$

\paragraph{Safe permutation scope.}
Pixel permutation gives identical non-DC one-row marginal laws and the
variance above, but not joint temporal exchangeability.  Uniform
acquisition-order permutation gives finite-population exchangeability
over its randomization draw.  If complementary masks are kept adjacent,
the exchangeable units are pair blocks, not individual frames.  Once a
single permutation is conditioned upon, the sequence is deterministic;
finite-population or randomization concentration should be invoked
rather than claiming that every realized order is stationary or mixing.

% End of paste-ready identifiability/lower-bound/Fisher/permutation material.
